\NativeTitle{57}{第七单元素养评价（A）}{（小数加减法）}{60}{4}{一 & 二 & 三 & 四}{ &  &  & }
\NativeSlot{278}{36}{%
\ExamGroup{一、填空。{\normalfont\kaishu\fontsize{10.8bp}{14bp}\selectfont （每空1分，共27分）}}

% p057_s01

\ExamQuestion{1.}{5}%
计算小数加减法时，要先把\mbox{（\hspace{1.8em}）}对齐，也就是把\mbox{（\hspace{1.8em}）}对齐。\par
}
\NativeSlot{331}{36}{%
% p057_s02

\ExamQuestion{2.}{5}%
填空：\(6.3+4.9+13.7=4.9+(\square+\square)\)；\allowbreak{}\(24.8-5.6-4.4=24.8-(\square\mathbin{\bigcirc}\square)\)。\par
}
\NativeSlot{387}{36}{%
% p057_s03

\ExamQuestion{3.}{5}%
两个加数的和是\(36.8\)，一个加数增加\(5.9\)，现在的和是\mbox{（\hspace{1.8em}）}。\par
}
\NativeSlot{412}{36}{%
% p057_s04

\ExamQuestion{4.}{5}%
20个\(0.1\)和20个\(0.01\)，它们的和是\mbox{（\hspace{1.8em}）}，差是\mbox{（\hspace{1.8em}）}。\par
}
\NativeSlot{437}{57}{%
% p057_s05

\ExamQuestion{5.}{5}%
按规律填数：\par\noindent （1）\(3.2,4.0,4.8,(\quad),(\quad)\)。\par

% p057_s06

\ExamQuestion{（2）}{5}%
\(10.1,8.9,7.7,(\quad),(\quad)\)。\par

% p057_s07

\ExamQuestion{（3）}{5}%
\(3.5,9.5,15.5,(\quad),27.5\)。\par
}
\NativeSlot{537}{36}{%
% p057_s08

\ExamQuestion{6.}{5}%
\(0.46\)扩大到原来的10倍，比原数大\mbox{（\hspace{1.8em}）}。\par
}
\NativeSlot{562}{36}{%
% p057_s09

\ExamQuestion{7.}{5}%
甲、乙两数的和是\(26.3\)，甲数是\(8.6\)，乙数是\mbox{（\hspace{1.8em}）}。\par
}
\NativeSlot{588}{36}{%
% p057_s10

\ExamQuestion{8.}{5}%
甲数是\(9.2\)，比乙数多\(3.7\)，乙数是\mbox{（\hspace{1.8em}）}。\par
}
\NativeSlot{612}{36}{%
% p057_s11

\ExamQuestion{9.}{5}%
\(4.8\)的计数单位是\mbox{（\hspace{4.5em}）}，它有\mbox{（\hspace{1.8em}）}个这样的计数单位，再添\mbox{（\hspace{1.8em}）}个这样的计数单位就是5。\par
}
\NativeSlot{664}{48}{%
\ExamQuestion{10.}{5}%
在圆圈里填上“＞”“＜”或“＝”。\par\smallskip
\noindent%
% p057_s12
\makebox[.499\linewidth][l]{\(5.6+4.4\ \bigcirc\ 12-1.7\)}%
% p057_s13
\makebox[.499\linewidth][l]{\(1.08+5.99\ \bigcirc\ 7.08\)}\par
\noindent%
% p057_s14
\makebox[.499\linewidth][l]{\(13.5-6\ \bigcirc\ 4.8+1.9\)}%
% p057_s15
\makebox[.499\linewidth][l]{\(7.3+5.9\ \bigcirc\ 8.2+4.7\)}\par
}
\NativeSlot{748}{36}{%
% p057_s16

\ExamQuestion{11.}{5}%
一根木条长\(18.2\)厘米，两段拼接重叠\(3.6\)厘米，拼好总长29厘米，另一根木条长\mbox{（\hspace{1.8em}）}厘米。\par
}
